\documentclass[11pt,letterpaper]{article}
\usepackage[top=0.55in,bottom=0.45in,left=0.65in,right=0.65in]{geometry}
\raggedbottom
\usepackage{times}
\usepackage{booktabs}
\usepackage{enumitem}
\usepackage[hidelinks]{hyperref}
\usepackage{titlesec}
\titlespacing*{\section}{0pt}{3pt}{1pt}
\titleformat{\section}{\normalfont\large\bfseries}{\thesection.}{0.5em}{}
\setlist[enumerate]{leftmargin=1.3em,topsep=1pt,itemsep=1pt,parsep=0pt}
\pagestyle{empty}

\title{\vspace{-3em}\textbf{Course Project Proposal}\\[1pt]
\normalsize SemEval-2026 Task 4: Narrative Story Similarity and Narrative Representation Learning}
\author{Claude Kwizera \\ \small CS 421: Natural Language Processing}
\date{\small October 16, 2026}

\begin{document}
\maketitle
\thispagestyle{empty}
\vspace{-3em}

\section{Task Selection}

This project enters SemEval-2026 Task 4 \cite{hatzel2026}, Track A.

The task gives the system three story summaries, an anchor and two candidates, and asks which candidate is
narratively closer to the anchor. The guidelines define ``closer'' along three aspects: course of
action, outcomes, and abstract theme. Annotators record which aspect drove each decision, and those records
become the diagnostic in Section~\ref{sec:approach}.

Only Track A is in scope. Track B asks for one vector per story, and the organisers compare those vectors
afterwards on triples the system never sees, so nothing can be encoded knowing which stories it will meet. That
rewards a general-purpose embedding. Track A shows both candidates at once, which rewards a comparator, and the
approach in Section~\ref{sec:approach} scores one story's cast against another's, so it needs both stories in hand.

How the data was built explains most of what follows. Summaries come from \emph{Tell Me Again}
\cite{tellmeagain}, cut to four to eight sentences, with a GPT-4o-mini filter that drops texts giving only a
premise. Each anchor is then sampled at random, and both candidates are drawn relative to it under
\texttt{story-emb} \cite{storyemb}, an encoder trained so that different summaries of one story land close
together: the nearest story is the assumed answer and the distractor comes from the 98th to 99th similarity
percentile, so both candidates sit close to the anchor. Two commercial LLMs then judge each triple, and only the
triples they disagree on are annotated. Of 1,237 such triples, 831 had two annotators agree and a third
annotation settled 208 more; the rest were dropped, leaving 1,039. Shallow cues are correspondingly weak: a
Jaccard baseline scores 45.0\% on dev, below the 50\% chance line, and 58.0\% on test, and TF--IDF cosine
scores 49.0\% and 58.25\%, against 78.00 for the best submitted system.

What might close that gap is structure: who does what to whom, which two retellings of one plot keep even when
they share few words. Prior work divides on which kind of structure. Event structure has not yet produced gains
here. The organisers trained narrative chain embeddings \cite{hatzel2023,chambers2008}, reached state of the art
on narrative cloze, the task of predicting a held-out event, and then reported that the representations ``perform
poorly at identifying narratively similar texts.'' In the shared task, StoryNet's symbolic narrative graphs
finished 43rd of 47 runs \cite{storynet}, and the winning team tried Propp's functions, a fixed inventory of
folktale plot roles, three times before dropping them \cite{cognac}. A related negative result holds more
generally: adding syntactic structure to a strong pretrained encoder has a ``limited and inconsistent'' effect on
language understanding tasks \cite{glavas2021}.

Character structure has produced gains, though on neighbouring tasks rather than on this one. Chaturvedi et al.\
\cite{chaturvedi2018} matched movie remakes. Their event-based model scored worse than bag-of-words on its own;
what helped was matching characters across the two stories, using agent and patient features from the dependency
parse, which raised precision at rank 1 from 0.558 to 0.637. They never resolved coreference, and they name
missed coreference as a source of error. Barancov\'a and Sileno \cite{barancova2026} add coreference-derived
character roles to dense embeddings and gain 14.6\% relative in retrieval accuracy over \emph{Tell Me Again}
($p<0.001$, four encoders). Neither tested parse-derived roles on triples filtered to defeat model agreement.
Whether that signal transfers to this harder setting is the hypothesis under test, and the project is built to
answer it either way.

\section{Official Metric}

Accuracy on the binary choice, and nothing else. The organisers report a human upper bound near 89\%
\cite{hatzel2026}, so the winning 78.00\% falls 11 points short of people rather than 22 points short of perfect.
They also report Krippendorff's $\alpha$ of 0.33 overall and 0.05 to 0.11 per aspect. Low agreement caps what any
system can reach and makes the per-aspect strata noisy, and 400 test items make a two-point difference unreliable
on its own. Both push Section~\ref{sec:approach} towards paired tests.

\section{Official Results}

\begin{center}
\small\begin{tabular}{lcc}
\toprule
& \textbf{Track A} & \textbf{Track B} \\
\midrule
Best system (COGNAC) & \textbf{78.00} & \textbf{72.00} \\
Strongest baseline & 67.00 \scriptsize GPT-4o-mini & 66.75 \scriptsize \texttt{e5} \\
Range, mean (ranked systems) & 53.50--78.00, 66.51 (44) & 54.50--72.00, 63.83 (27) \\
Human upper bound & \multicolumn{2}{c}{$\approx 89$} \\
\bottomrule
\end{tabular}
\end{center}

Forty-six teams sent 71 submissions \cite{hatzel2026}. Figures above are from the overview and system papers
\cite{hatzel2026,cognac,e5}; the two means are computed here, and the range covers ranked systems only,
so the baselines sit outside it. The strongest Track A baseline is a single zero-shot GPT-4o-mini prompt at
67.00, about 22nd of 44: beating it is easy, beating the field is not.

\section{Data Access}

The data downloaded with no trouble: public on GitHub, no sign-up and no licence form
(\url{https://github.com/narrative-similarity-task/narrative-similarity-dataset}). Track A has 39 sample, 200 dev
and 400 test triples; Track B has 849 test stories. Summaries and labels are CC BY-SA 4.0, labels also CC0.

Three properties of the release shape the plan. First, no training split is provided at all, so there is no
fine-tuning: the approach is feature engineering with cross-validation on the labelled splits. Second, the gold
test labels are public, which allows rebuilding the leaderboard offline without CodaBench but makes overfitting
easy; one team was disqualified for exactly that \cite{hatzel2026}. I have run the lexical baselines on test
once, to check the scorer against the published figures; from here the test set is closed until a single final
run. Third, the
splits share stories even where they share no triple: the Track A and Track B test sets have 265 stories in
common, and 190 of the 848 Track A test stories appear in dev or sample. Features tuned on dev
can therefore see test stories, so L5 is reported in dev-tuned and dev-free variants.

\section{Planned Approach}\label{sec:approach}

Each layer adds one representation and tests one step of the argument above. \textbf{L0} is random and
Jaccard, which reproduce the published baselines; this implementation scores 58.0\% on test against the published
56.25\%, a seven-item difference attributable to tokenisation and to be reconciled first.
\textbf{L1} is TF--IDF cosine, which measures how much of the task is word overlap. \textbf{L2} is sentence
embeddings (\texttt{e5} \cite{e5}); the organisers publish no Track A figure for it, so L2 establishes one.
\textbf{L3} adds characters: coreference chains from \texttt{maverick-coref} \cite{maverick,jm3} give each story
its cast, and story similarity is the score of the best one-to-one matching between the two casts, computed with
the Hungarian algorithm over mention-string and context similarity. \textbf{L4} adds roles: for each character, the verb
lemmas where it is the syntactic subject act as an agent proxy and those where it is the object as a patient
proxy \cite[Chapter~22]{jm3}, and the matching is scored on shared role--lemma pairs instead of mentions. \textbf{L5}
combines L2 with the L3 and L4 scores, weights fitted by cross-validation. Each layer scores the anchor against
both candidates and picks the higher.

No maintained PropBank SRL system exists (AllenNLP was archived in 2022, and SRL by LLM prompt reaches 48.6\%
argument match against 88.3\% for humans \cite{bonn2024}), so L4 uses dependency proxies, as Chaturvedi et al.\
did \cite{chaturvedi2018}.

A first attempt at the per-aspect diagnostic is why the statistics below are strict. Splitting test by
course of action, TF--IDF scores 62.4\% on the 186 items annotators did not mark and 54.7\% on the 214 they did.
The same split on dev reverses, 44.3\% against 53.4\%, and neither gap is significant ($z=1.56$, $p=0.12$;
$z=-1.28$, $p=0.20$). The instability has an identifiable cause: with two annotators, ``did not mark'' pools 31
test items that neither annotator marked with 136 where one of the two did. The diagnostic will therefore use
unanimous strata only, and will use outcomes rather than course of action, where the clean strata hold 79 and 133
items rather than 188 and 32. Abstract theme is unusable for this purpose: every annotator marked it on 392 of
the 400 test items.

Evaluation is official accuracy against the gold labels. Bootstrap confidence intervals \cite[Section~4.11.1]{jm3}
describe a single system; two systems are compared with McNemar's test on the same items \cite{dietterich},
since at these sample sizes overlapping intervals would hide real differences. Model selection
is cross-validation over dev plus sample (239 items), with configurations capped in advance.

L0 to L2 need no training and L3 and L4 only run pretrained models, so Colab Pro is enough and no paid API is needed
before November 25. The project is in English.

\section{Failure Cases and Mitigations}

\begin{enumerate}[label=\textbf{F\arabic*.}]
\item The character and role features add nothing over the embeddings. This is the most likely outcome, and there
is precedent \cite{hatzel2023,glavas2021,chaturvedi2018}. The ablation makes that a reportable result, and the
diagnostic says in which stratum the signal is absent.
\item Coreference is unreliable on summaries this short, and coreference scores have been shown not to predict
downstream gains in at least one other domain \cite{chai2022}. I will check 50 items by hand early, report the
error rate, and use the singleton-keeping LitBank model \cite{maverick}.
\item A gain turns out to be noise, as the stratified result above already was. Guards: McNemar's test,
bootstrap intervals on both splits, capped configurations, one clean test run.
\item Recent work overlaps \cite{barancova2026}: it ranks unfiltered summaries with LLM-assigned roles, where
this project classifies filtered triples with parse-derived roles.
\end{enumerate}

\newpage
\begin{thebibliography}{99}
\footnotesize\setlength{\itemsep}{0pt}
\bibitem{hatzel2026} Hans Ole Hatzel, Ekaterina Artemova, Haimo Stiemer, Evelyn Gius, and Chris Biemann. SemEval-2026 Task 4: Narrative Story Similarity and Narrative Representation Learning. \emph{SemEval-2026}, 3460--3478. \url{https://aclanthology.org/2026.semeval-1.429/}
\bibitem{cognac} Tisa Islam Erana, Azwad Anjum Islam, Anshu Kiran Sharma, and Mark A. Finlayson. COGNAC at SemEval-2026 Task 4: Evaluating Narrative Components with LLMs for Hard Story Similarity Cases. \emph{SemEval-2026}, 2290--2300. \url{https://aclanthology.org/2026.semeval-1.290/}
\bibitem{storynet} Amol Vijayachandran, Ananth Rajesh, Siddharth Mago, Maitreya Prafulla Chitale, and Aparajitha Allamraju. IIITH Boys at SemEval-2026 Task 4: StoryNet---Understanding Narrative Story Similarity through Symbolic Representations. \emph{SemEval-2026}, 3382--3393. \url{https://aclanthology.org/2026.semeval-1.420/}
\bibitem{chaturvedi2018} Snigdha Chaturvedi, Shashank Srivastava, and Dan Roth. Where Have I Heard This Story Before? Identifying Narrative Similarity in Movie Remakes. \emph{NAACL-HLT 2018}, 673--678. \url{https://aclanthology.org/N18-2106/}
\bibitem{hatzel2023} Hans Ole Hatzel and Chris Biemann. Narrative Cloze as a Training Objective: Towards Modeling Stories Using Narrative Chain Embeddings. \emph{5th Workshop on Narrative Understanding}, 118--127. \url{https://aclanthology.org/2023.wnu-1.19/}
\bibitem{barancova2026} Alexandra Barancov\'a and Giovanni Sileno. Narrative Similarity Detection Using Character Roles. \emph{Text2Story'26 Workshop}, CEUR-WS Vol-4202. \url{https://ceur-ws.org/Vol-4202/paper6.pdf}
\bibitem{chambers2008} Nathanael Chambers and Dan Jurafsky. Unsupervised Learning of Narrative Event Chains. \emph{ACL-08: HLT}, 789--797. \url{https://aclanthology.org/P08-1090/}
\bibitem{tellmeagain} Hans Ole Hatzel and Chris Biemann. Tell Me Again! a Large-Scale Dataset of Multiple Summaries for the Same Story. \emph{LREC-COLING 2024}, 15732--15741. \url{https://aclanthology.org/2024.lrec-main.1366/}
\bibitem{storyemb} Hans Ole Hatzel and Chris Biemann. Story Embeddings---Narrative-Focused Representations of Fictional Stories. \emph{EMNLP 2024}, 5931--5943. \url{https://aclanthology.org/2024.emnlp-main.339/}
\bibitem{e5} Liang Wang, Nan Yang, Xiaolong Huang, Binxing Jiao, Linjun Yang, Daxin Jiang, Rangan Majumder, and Furu Wei. Text Embeddings by Weakly-Supervised Contrastive Pre-training. arXiv:2212.03533, 2022.
\bibitem{maverick} Giuliano Martinelli, Edoardo Barba, and Roberto Navigli. Maverick: Efficient and Accurate Coreference Resolution Defying Recent Trends. \emph{ACL 2024}, 13380--13394. \url{https://aclanthology.org/2024.acl-long.722/}
\bibitem{glavas2021} Goran Glava\v{s} and Ivan Vuli\'c. Is Supervised Syntactic Parsing Beneficial for Language Understanding Tasks? An Empirical Investigation. \emph{EACL 2021}, 3090--3104. \url{https://aclanthology.org/2021.eacl-main.270/}
\bibitem{chai2022} Haixia Chai, Nafise Sadat Moosavi, Iryna Gurevych, and Michael Strube. Evaluating Coreference Resolvers on Community-based Question Answering: From Rule-based to State of the Art. \emph{CRAC 2022}, 61--73. \url{https://aclanthology.org/2022.crac-1.7/}
\bibitem{bonn2024} Julia Bonn, Harish Tayyar Madabushi, Jena D. Hwang, and Claire Bonial. Adjudicating LLMs as PropBank Annotators. \emph{DMR 2024 @ LREC-COLING}, 112--123. \url{https://aclanthology.org/2024.dmr-1.12/}
\bibitem{jm3} Daniel Jurafsky and James H. Martin. \emph{Speech and Language Processing}, 3rd ed. Online manuscript released August 19, 2026. \url{https://web.stanford.edu/~jurafsky/slp3/}
\bibitem{dietterich} Thomas G. Dietterich. Approximate Statistical Tests for Comparing Supervised Classification Learning Algorithms. \emph{Neural Computation} 10(7):1895--1923, 1998.
\end{thebibliography}

\end{document}
